\documentclass[12pt,letterpaper]{article}

%Packages
\usepackage{amsmath}
\usepackage{amsthm}
\usepackage{amsfonts}
\usepackage{amssymb}
\usepackage{amscd}
\usepackage{dsfont}
\usepackage{enumerate}
\usepackage{fancyhdr}
\usepackage{mathrsfs}
\usepackage{bbm}
\usepackage{framed}
\usepackage{mdframed}
\usepackage{cancel}
\usepackage{float}
\usepackage{mathtools}
%\usepackage[]{mcode}
\usepackage{graphicx}
\usepackage{tikz}
\usepackage{hyperref}
\hypersetup{
    colorlinks=true,
    linkcolor=blue,
    filecolor=magenta,      
    urlcolor=cyan,
}
\usetikzlibrary{automata,arrows}

%Page formatting
\usepackage[letterpaper,voffset=-.5in,bmargin=4cm,footskip=1cm]{geometry}
\setlength{\parindent}{0.0in}
\setlength{\parskip}{0.1in}
\allowdisplaybreaks
\headheight 15pt
\headsep 10pt

% Common Commands
\newcommand\N{\mathbb N}
\newcommand\Z{\mathbb Z}
\newcommand\R{\mathbb R}
\newcommand\Q{\mathbb Q}
\newcommand\lcm{\operatorname{lcm}}
\newcommand\setbuilder[2]{\ensuremath{\left\{#1\;\middle|\;#2\right\}}}
\newcommand\E{\operatorname{E}}
\newcommand\V{\operatorname{V}}
\newcommand\Pow{\ensuremath{\operatorname{\mathcal{P}}}}

\DeclarePairedDelimiter\ceil{\lceil}{\rceil}
\DeclarePairedDelimiter\floor{\lfloor}{\rfloor}

\newcommand\definition{\textbf{Definition: }}
\newcommand\proposition{\textbf{Proposition: }}

% 22 style
\newcommand\hint[1]{\textbf{Hint}: #1}
\newcommand\note[1]{\textbf{Note}: #1}
\newenvironment{22enumerate}{\begin{enumerate}[a.]\itemsep0em}{\end{enumerate}}
\newenvironment{22itemize}{\begin{itemize}\itemsep0em}{\end{itemize}}
\fancypagestyle{firstpagestyle} {
  \renewcommand{\headrulewidth}{0pt}%
  \lhead{\textbf{CSCI 0220}}%
  \chead{\textbf{Discrete Structures and Probability}}%
  \rhead{Littman}%
}



\pagestyle{fancyplain}

\newcommand\EX{\mathds{E}}
\newcommand\PR{\mathds{P}}
\newcommand\zlam{Z_\lambda}
\DeclareMathOperator*{\argmin}{arg\,min}

%%%%%%%%%%%%%%%% COMPILE WITH \soltrue or \solfalse %%%%%%%%%%%%%%%%%%%%%%%%%%%%%
\newif\ifsol
\soltrue  % SOLUTIONS
% \solfalse % NO SOLUTIONS
%%%%%%%%%%%%%%%%%%%%%%%%%%%%%%%%%%%%%%%%%%%%%%%%%%%%%%%%%%%%%%

%%%%%%%%%%%%%%%% solution help %%%%%%%%%%%%%%%%%%%%%
\newcommand{\solu}[2]{ \begin{mdframed} \ifsol #2 \else \vspace{#1} \fi \end{mdframed} }
\newcommand{\solt}{ \ifsol (T) \fi }
\newcommand{\solf}{ \ifsol (F) \fi }
\newcommand{\solm}[1]{\ifsol \textit{(#1)} \fi}


\begin{document}
  \thispagestyle{firstpagestyle}
  \begin{center}
    {\large \textbf{Recitation 3}}
    
    {\large Set Equivalence}

  \end{center}

\section*{Set Theory: Review}

Recall the definitions and terminology for set theory:

\textbf{Defn 1:} A \textbf{set} is a collection of objects with no repetition or order.

\textbf{Defn 2:} $B$ is a \textbf{subset} of $A$ if every element in $B$ is also in $A$. This is written as $B \subseteq A$.

\textbf{Defn 3:} $\Pow{(A)}$, called the \textbf{power set} of $A$, is the set of all subsets of $A$.

\textbf{Defn 4:} $A \cup B$ denotes the \textbf{union} of sets $A$ and $B$. This contains all the elements from $A$, and all of the elements from $B$. 
	
\textbf{Defn 5:} $A \cap B$ denotes the \textbf{intersection} of sets $A$ and $B$. This contains only the elements that appear in both of the sets.
	
\textbf{Defn 6:} $A - B$ denotes the \textbf{difference} of sets $A$ and $B$. This contains elements that appear in $A$ but not $B$. In other words, $A - B = \{x \in A \ | \ x \notin B\}$.
	
\textbf{Defn 7:} $\overline{A}$ denotes the \textbf{complement} of $A$ relative to some universal set $U$. $\bar{A} = U - A$, that is, it is everything except what is in $A$.
	
\textbf{Defn 8:} $|A|$ denotes the \textbf{cardinality} of $A$, which is a count of the number of elements contained in $A$.
\newpage
\textbf{Task 1}

Define the following sets
\begin{align*}
	A &= \left\{\varnothing, 0, \left\{\varnothing\right\},
\left\{0,\varnothing\right\} \right\} \\
	B &= \left\{\varnothing, \left\{\varnothing\right\},
\left\{0, \varnothing, 1\right\} \right\} \\
	C &= \{x \mid \exists \ y \in \Z \text{ s.t. } y^2=x, \ x < 10 \} \\
    D &= \{a, b, c, d, e\} \\
    E &= \{c, d, e\} \\
    F &= \{d, e, f, g\} \\
\end{align*}
\begin{enumerate}[a.]
    \item Find the following sets:
  		\begin{enumerate}[i.]
  			\item $A \cup B$
  			\begin{mdframed} \vspace{0.5cm} \end{mdframed}
      		\item $A \cap B$
      		\begin{mdframed} \vspace{0.5cm} \end{mdframed}
      		\item $A - B$
      		\begin{mdframed} \vspace{0.5cm} \end{mdframed}
      		\item $\Pow(B)$ ($=2^B$), the power set of $B$
      		\begin{mdframed} \vspace{0.5cm} \end{mdframed}
      		\item $C - (A \cup B)$
      		\begin{mdframed} \vspace{0.5cm} \end{mdframed}
      		\item $\{x \mid x \in B, |x| \notin C\}$
      		\begin{mdframed} \vspace{0.5cm} \end{mdframed}
            \item $A \times E$
            \begin{mdframed} \vspace{0.5cm} \end{mdframed}
            \item \textit{Optional: }$(D \times E) - (E \times F)$
            \begin{mdframed} \vspace{0.5cm} \end{mdframed}
            \item \textit{Optional: }$(D - E) \cap (D \cap E)$
            \begin{mdframed} \vspace{0.5cm} \end{mdframed}
	   \end{enumerate}
  \item Find the cardinalities of the following sets:
        \begin{enumerate}[i.]
          \item $C$
          \begin{mdframed} \vspace{0.5cm} \end{mdframed}
          \item $A \times B$
          \begin{mdframed} \vspace{0.5cm} \end{mdframed}
          \item $\Pow(A \cup B)$
          \begin{mdframed} \vspace{0.5cm} \end{mdframed}
          \item $F - (D \cap E)$
          \begin{mdframed} \vspace{0.5cm} \end{mdframed}
          \item \textit{Optional: }$\Pow(\Pow(\varnothing))$
          \begin{mdframed} \vspace{0.5cm} \end{mdframed}
        \end{enumerate}
\end{enumerate}

\textbf{Checkpoint 1 — Call over a TA!}

\newpage
\section*{Set Equivalence}
\subsection*{Set Element Method}
\textbf{Defn 1:} Two sets $S$ and $T$ are equal if and only if they have the same elements:
\begin{equation*}
    S = T \iff (\forall x : x \in S \iff x \in T)
\end{equation*}
\textbf{Defn 2:} S and T are equal if and only if both S is a subset of T and T is a subset of S.

From these two definitions, we can construct our set-element method for proving set equalities; that is, we show that two sets are equal if and only if every element of $S$ is also an element of $T$ and every element of $T$ is also an element of $S$.

In other words, to prove that one set is a subset of the other, we can show that $S \subseteq T$ for arbitrary sets $S$ and $T$ using the following steps:
\begin{enumerate}[1.]
    \item Let x be an element of S.
    \item Prove that x is an element of T.
    \item Conclude that $S \subseteq T$.
\end{enumerate}

\textbf{Example}

\textit{Claim: } $A \cap (A \cup B) = A$.

\textit{Proof: } We show that both $A \cap (A \cup B) \subseteq A$ and $A \subseteq A \cap (A \cup B)$.
\begin{enumerate}
    \item We will first prove our sub-claim that $A \cap (A \cup B) \subseteq A$.
    
    Consider any $x \in A \cap (A \cup B)$. This means that $x \in A$ and $x \in A \cup B$. In particular, we know that $x \in A$. Thus, we have shown that $A \cap (A \cup B) \subseteq A$.
    
    \item We will then prove the second sub-claim that $A \subseteq A \cap (A \cup B)$.
    
    Consider any $x \in A$. Then, $x \in A \cup B$. Thus, we know that $x \in A$ and $x \in A \cup B$, which can be rewritten as $x \in A \cap (A \cup B$. This shows that $A \subseteq A \cap (A \cup B)$.
\end{enumerate}
Therefore, $A = A \cap (A \cup B)$.

Now, let's practice!

\newpage
\textbf{Task 2}

For each of these statements, either prove (using the 'set element' method) or disprove the claim.

Hint: Use a counterexample to disprove a claim!
\begin{enumerate}[a.]
    \item For any two sets $A$ and $B$, $\Pow(A \cap B) = \Pow(A) \cap \Pow(B)$. 
    \begin{mdframed} \vspace{4cm} \end{mdframed}
        

    \item For any two sets $A$ and $B$, $\Pow(A \cup B) = \Pow(A) \cup \Pow(B)$.
    \begin{mdframed} \vspace{4cm} \end{mdframed}
    
    \end{enumerate}

\newpage
\textbf{Task 3}

\begin{enumerate}[a.]    
    \item Let $A$ and $B$ be subsets of some universal set. Then $A - (A - B) = A \cap B$.
    \begin{mdframed} \vspace{4cm} \end{mdframed}
    
    \item \textit{Optional:} Let $A$ and $B$ be the following sets:
    \begin{equation*}
    \begin{aligned}
        A =& \{6n : n \in \mathbb{Z}\} \\
        B =& \{2n : n \in \mathbb{Z}\} \cap \{3n :  n \in \mathbb{Z}\}
    \end{aligned}
    \end{equation*}
    Prove that $A = B$.
    \begin{mdframed} \vspace{4cm} \end{mdframed}
    
    \item \textit{Optional:} Prove that $A \times (B \cup C) = (A \times B) \cup (A \times C)$.   
    \begin{mdframed} \vspace{4cm} \end{mdframed}
    
\end{enumerate}

\textbf{Checkpoint 2 — Call over a TA!}

\newpage
\subsection*{Set Algebra}
We can also prove that two sets are equivalent using set algebra. More concretely, we can use the stated laws of set algebra to convert one side of the equation to the other (or convert both sides to an identical expression).

\textbf{Example} (from the \href{https://cs22.netlify.app/#resources}{sample proofs} on our website):

Prove that $(A \cap B) \cup (A - B) = A \cap (B \cup (A - B))$.

\begin{flalign*}
& (A \cap B) \cup (A - B) \\
&= (A \cap B) \cup (A \cap \overline{B}) && \text{(Set Difference Law)} \\
&= (A \cap (B \cup \overline{B}) && \text{(Distribution)} \\
&= A \cap U && \text{(Complement Law)} \\
&= A && \text{(Identity Law)} \\
&= A \cap (A \cup B) && \text{(Absorption)} \\
&= A \cap (B \cup A) && \text{(Commutivity)} \\
&= A \cap ((B \cup A) \cap U) && \text{(Identity Law)} \\
&= A \cap ((B \cup A) \cap (B \cup \overline{B}) && \text{(Complement Law)} \\
&= A \cap (B \cup (A \cap \overline{B})) && \text{(Distribution)} \\
&= A \cap (B \cup (A - B)) && \text{(Set Difference Law)}
\end{flalign*}
Therefore, the equality holds since all these steps are biconditionally true.

\newpage
\textbf{Task 4}
\begin{enumerate}[a.]
    \item Let $A$ and $B$ be subsets of some universal set $U$. Prove that $(A \cap \overline{B}) \cup B = A \cup B$ using set algebra.
    \begin{mdframed} \vspace{4cm} \end{mdframed}

    \item Let $A$ and $B$ be subsets of some universal set $U$. Prove that $(A - B) \cup (B - A) = (A \cup B) - (A \cap B)$ using set algebra.
    
    \begin{mdframed} \vspace{4cm} \end{mdframed}

    \item \textit{Optional: }Show that $(A - B) - (B - C) = A - B$.
    \begin{mdframed} \vspace{4cm} \end{mdframed}

\end{enumerate}

\newpage
\section*{Google Ad Personalization}
\textbf{Task 5}
 \begin{enumerate}[a.]
    \item Read \href{https://www.practicalecommerce.com/how-google-collects-data-to-personalize-ads}{this article} on Google's ad personalization, and how it collects data on its users to select targeted ads. 
    
    \item If you have a Google account that you actively use, follow the instructions on \href{https://www.businessinsider.com/what-does-google-know-about-me-search-history-delete-2019-10#first-navigate-to-your-google-account-homepage-by-clicking-the-widget-in-the-top-right-corner-of-any-google-site-1}{this page} and look at the groups that Google believes you are a part of. If you have ad penalization enabled and feel comfortable doing so, skim your entire ad profile. Did anything surprise you? Why or why not?
    
    \begin{mdframed} \vspace{2cm} \end{mdframed}
    
    \item 
    \textbf Google defines your ad profile as the intersection of several traits, or sets. This profile will determine what ads you are shown - advertisers choose what demographic groups and interest groups they want their ads to target (you can learn more about interest targeting \href{https://support.google.com/google-ads/answer/2497941}{here} and demographic targeting \href{https://support.google.com/google-ads/answer/2580383}{here}). Think about the assumptions and biases that come with putting people into sets. What problems and inaccuracies could arise by defining what content individuals receive solely based on the sets they belong to?
    
    \begin{mdframed} \vspace{4cm} \end{mdframed}
    
    \item Being at the intersection of several traits means that people's experiences could differ drastically from people in one set or the other. For example, take two traits that tend to be underprivileged: low education level and low income level. How might a upper-class high school dropout's opportunities and livelihood differ from a high school dropout experiencing poverty? What are some other cases of people who belong to the same set in one instance but experience vastly different amounts of privilege and discrimination based on what other sets they are a member of?
    
    % Is describing their experience as simple as combining the experience of one group and the experience of the other group, or is there something unique about being at the intersection of several groups?
    % What are other cases where people who both identify with one set might experience privilege and discrimination differently because of the existence of intersections?
    
    \begin{mdframed} \vspace{4cm} \end{mdframed}
    
    \item When it comes to data analysis and storage, it is common practice to group people and profiles into discrete sets such as race, gender identity, and geographical location. Standard database design requires people to be grouped together based on certain traits, with each trait being represented by a column. What potential conflicts and problems could occur when trying to fit people into a finite number of sets? How could assumptions that apply to the entire set harm the individuals in them?
    
    \begin{mdframed} \vspace{4cm} \end{mdframed}
\end{enumerate}
\subsection*{Checkoff — Call over a TA!}

\end{document}
